\documentclass[11pt]{article}
\usepackage{fontspec,amsmath,amssymb,amsthm,hyperref}
\usepackage[margin=1in]{geometry}
\setlength{\emergencystretch}{3em}
\title{Quadratic integrability and exact energy of the synchronous process}
\author{LangevinLean mathematical exposition}
\date{}
\begin{document}\maketitle
\noindent Source module: \texttt{AverageLMC/PathwiseEnergy.lean}. Compile this proof document with XeLaTeX or LuaLaTeX. The accompanying HTML pairs the mathematical argument with the exact Lean statements and records the scope of the formalization.


All expectations and derivatives in the stochastic proof are justified on the actual space of square-integrable random variables. A smooth potential need not be globally Lipschitz: its quadratic growth is the correct condition here. The common additive noise cancels before the squared error is differentiated.



Writing \(\langle U,V\rangle_{L^2}=\mathbb E\langle U,V\rangle\), strong convexity and the common-noise integral equation give


\[2\langle X-Y,\nabla f(Y)\rangle_{L^2}\le 2\bigl(\mathbb Ef(X)-\mathbb Ef(Y)\bigr)-m\|X-Y\|_{L^2}^{2}.\]
\[\|X+B_h-Y_h\|_{L^2}^{2}=\|X-Y_0\|_{L^2}^{2}+2\int_0^h\langle X+B_t-Y_t,\nabla f(Y_t)\rangle_{L^2}\,dt.\]
\section{Smooth potentials are integrable on L\ensuremath{{}^{2}}}


Apply the quadratic upper model both to f and to −f at zero. These two bounds control the absolute value of f(x) by the sum of |f(0)|, the product of the initial gradient norm and ‖x‖, and L‖x‖\ensuremath{{}^{2}}/2. A square-integrable random vector has integrable norm and squared norm, so comparison proves integrability of f evaluated at its L\ensuremath{{}^{2}} representative. The second theorem transfers the result from the quotient-space representative to any supplied random vector with finite second moment by its almost-everywhere equality with MemLp.toLp.



\noindent Lean declarations: \texttt{integrable\_smooth\_compLp}, \texttt{integrable\_smooth\_comp}.

\section{Integrate the strong supporting plane}


At each outcome, strong convexity gives f(X)−f(Y) at least the inner product of ∇f(Y) with X−Y plus m‖X−Y‖\ensuremath{{}^{2}}/2. Every scalar term is integrable by quadratic growth, and the inner product of two L\ensuremath{{}^{2}} vectors is integrable by the L\ensuremath{{}^{2}} Cauchy–Schwarz inequality. Integration, the definition of the L\ensuremath{{}^{2}} inner product, and the exact squared-norm integral identity give the supporting-plane inequality on L\ensuremath{{}^{2}} itself. The lifted vector field is the actual gradient applied pointwise, rather than a supplied random drift.



\noindent Lean declarations: \texttt{lp\_strong\_gradient\_support}.

\section{Differentiate after the common noise cancels}


For a continuous vector field g, the primitive A+∫₀ᵗg has derivative g(t). The derivative of its squared Hilbert norm is twice its inner product with g(t); the fundamental theorem of calculus proves the displayed exact energy identity. In the Langevin integral equation, X+B(t)−Y(t) equals X−Y(0)+∫₀ᵗ∇f(Y(s))ds because the same B drives both processes. Substitute this expression into the Hilbert identity. No derivative of a Brownian path and no stochastic isometry is assumed.



\noindent Lean declarations: \texttt{integral\_curve\_norm\_sq}, \texttt{langevin\_energy\_identity}.

\end{document}
